% Generated by scripts/assemble.py; corrections belong in source/corrections.json.
\examhead{145}{中南大学 2025 年数学分析试卷}{138}
% p139-b06; source PDF page 139
\par\Needspace{5\baselineskip}
\blockanchor{p139-b06}
\majorquestion{一}计算题. 每题 10 分, 共 40 分.\par
% p139-b07; source PDF page 139
\blockanchor{p139-b07}
(1) $\lim_{x\to 0^{+}}\frac{\sin(\tan x) - \tan(\sin x)}{\sin^4 x}.$\par
% p139-b08; source PDF page 139
\blockanchor{p139-b08}
(2) $\int_0^2\frac{\mathrm{d}x}{e^x + e^{-x}}.$\par
% p139-b09; source PDF page 139
\blockanchor{p139-b09}
(3) $\lim_{n\to \infty}\left(1 - \frac{1}{n} +\frac{1}{n^2}\right)^{1 - n}.$\par
% p139-b10; source PDF page 139
\blockanchor{p139-b10}
(4) $\int_0^1\frac{\cos x\cos(\sin x)}{\sin(\sin x) + \cos(\sin x)}\mathrm{d}x.$\par
% p139-b11; source PDF page 139
\par\Needspace{5\baselineskip}
\blockanchor{p139-b11}
\majorquestion{二}(15 分) 确定如下级数 $\sum_{n=0}^{\infty} \frac{2n+1}{3^{n+2}} x^{2n}$ 的收敛域, 并计算在收敛域内的和函数.\par
% p139-b12; source PDF page 139
\par\Needspace{5\baselineskip}
\blockanchor{p139-b12}
\majorquestion{三}(15 分) 证明: 函数 $y = \frac{x}{1 + x^4}$ 在实数域 $\mathbb{R}$ 上有界且一致连续.\par
% p139-b13; source PDF page 139
\par\Needspace{5\baselineskip}
\blockanchor{p139-b13}
\majorquestion{四}(15 分) 计算曲线 $x = e^{t} \cos t, y = e^{t} \sin t$ 在 $t \in [0, \pi]$ 的长度.\par
% p139-b14; source PDF page 139
\par\Needspace{5\baselineskip}
\blockanchor{p139-b14}
\majorquestion{五}(15 分) 求由方程 $x^{2} + 2xy + 2y^{2} = 1$ 所确定的隐函数的极值.\par
% p139-b15; source PDF page 139
\par\Needspace{5\baselineskip}
\blockanchor{p139-b15}
\majorquestion{六}(15 分) 证明: 若函数 $u = f(x, y)$ 满足拉普拉斯方程 $u_{xx} + u_{yy} = 0$, 则函数 $v = f(x - \sqrt{3}y, \sqrt{3}x + y)$ 也满足拉普拉斯方程.\par
% p139-b16; source PDF page 139
\par\Needspace{5\baselineskip}
\par\Needspace{6\baselineskip}
\blockanchor{p139-b16}
\majorquestion{七}(15 分) 设 $f: (0,1) \to \mathbb{R}$ 满足以下条件 (*)：存在 $L > 0, \alpha \in (0,1)$，使得对任意的 $x, y \in (0,1)$，始终有\par
\nopagebreak[4]
% p139-b17; source PDF page 139
\blockanchor{p139-b17}
\displayformula{| f (x) - f (y) | \leq L | x - y | ^ {\alpha}.}
% p139-b18; source PDF page 139
\blockanchor{p139-b18}
(1) (6 分) 证明: $f$ 在 $(0,1)$ 上连续.\par
% p139-b19; source PDF page 139
\blockanchor{p139-b19}
(2) (6 分) 证明: $f$ 在 $(0,1)$ 上一致连续.\par
% p139-b20; source PDF page 139
\blockanchor{p139-b20}
(3) (3 分) 如果函数 $f$ 满足条件 (*), 则称其为 $C^\alpha$ 函数, 请举出一个例子, 满足对于 $0 < \beta < \alpha$, 该函数是 $C^\beta$ 函数, 但它不是 $C^\alpha$ 函数.\par
% p139-b21; source PDF page 139
\par\Needspace{5\baselineskip}
\blockanchor{p139-b21}
\majorquestion{八}(20 分) 解答如下问题:\par
% p139-b22; source PDF page 139
\blockanchor{p139-b22}
(1) (10 分) 计算 $f(x) = \frac{2x\sin\theta}{1 - 2x\cos\theta + x^2}$ 在 $x = 0$ 处的泰勒展开式.\par
% p139-b23; source PDF page 139
\blockanchor{p139-b23}
(2) (10 分) 设 $|x| < 1$, 计算 $\int_0^\pi \ln (1 - 2x\cos \theta + x^2)\mathrm{d}\theta$.\par
